Selecting products with small $R$ raises the mean numerical systolic
ratio, but further restriction within that selection need not raise it
again. In a fresh pool of 100,000 Lagrangian products, selecting the ten
lowest-$R$ bodies per side-count group gave mean ratio $0.626$, compared
with $0.316$ for the same number of comparison bodies. Selection used only geometry, before capacity was evaluated.
Each factor-side-count group contributed equally. The comparison bodies
came from the same pool, excluded the selected bodies, and were chosen
without evaluating their capacities.

The distribution of area among the ridges supplies a second comparison.
When the total symplectic ridge area is positive, define
\[
 C(K)=\frac{\max_F A_\omega(F)}{\sum_F A_\omega(F)}.
\]
This is the share of the total carried by the largest ridge; like $R$,
it is invariant under translation, common dilation and linear symplectic
maps. Among products in the
lowest one percent by $R$ in a separate fresh sample, the half with lower
$C$ had a mean advantage of $0.034$ over the other half, positive in eight
of ten groups. Thus largest-ridge dominance distinguishes mean ratios within this
selected pool, even when the total ridge area is small.

The strong ordering in the broad samples is not retained within the
low-$R$ panels themselves. Across 1,000 evaluated products in the lowest
one percent by $R$, correlations within the ten side-count groups range
from $-0.315$ to $+0.301$; pooling ranks computed within each group gives
$+0.032$. The most extreme tenth has a larger mean ratio than the rest
in only two groups. Small $R$ distinguishes this population from the
comparison bodies, but does not reliably order the ratios within it.
